\section{Projection and decoding statements}
For any linear code $C\leq\mathbb F_p^N$, the condition
$C+C^\perp=\mathbb F_p^N$ is equivalent to $C\cap C^\perp=0$, by the
dimension formula. This is source Lemma~1. The shared symmetric-idempotent
argument gives the actual orthogonal projections, rather than just a
rank test.

If $A(A+I)=0$, the projections onto $C$ and $RC$ are $-A$ and $I+A$.
Conversely $C\perp RC$ implies this identity, so containment already
implies $C=(RC)^\perp$. This recovers source Proposition~1(1),
Proposition~2(1) and quoted Result~2 from \Uref{ULCD.Reflexive}.
The unconditional inclusion in Result~2 is also immediate:
if $Av=0$, then $v=(A+I)v$, so $C^\perp\subseteq RC$.
For a received row vector $v=\sum_xv(x)e_x$, with adjacency and reflexive
rows $r_x,s_x$, respectively, the decomposition is
\[
 v=-\sum_xv(x)r_x+\sum_xv(x)s_x.
\]
In characteristic two, if $v$ already belongs to one summand, it is the
sum of its corresponding rows indexed by $\operatorname{Supp}(v)$.

Suppose the graph is regular of degree $\nu$ and $AB=0$.
Put $\bar\nu=N-1-\nu$, with these integers also denoting their field
images. \Uref{ULCD.Complement} states that both codes are LCD and gives
\[
\begin{array}{c|c|c}
 &\text{orthogonal decomposition}&\text{projection matrices}\\\hline
\nu\ne0&\mathbb F_p^N=C\perp\!\oplus\bar C&J-A,\ -B\\
\bar\nu\ne0&\mathbb F_p^N=C\perp\!\oplus\bar C&-A,\ J-B\\
\nu=\bar\nu=0&\mathbb F_p^N=C\perp\!\oplus\bar C\perp\!\oplus\langle\mathbf j\rangle
 &-A,\ -B,\ J.
\end{array}
\]
Thus the exceptional case in source Proposition~2(2) occurs exactly when
both degrees vanish. Its defect is exactly one and is the displayed
constant line; no integral-spectrum assumption is necessary. When the
codes are complementary this is precisely Proposition~1(2), including
its distinction according to whether $p\mid\nu$.

The relation between the two reflexive tests is
\begin{equation}\label{eq:two-reflexive}
 B(B+I)=A(A+I)+(N-2\nu-1)J.
\end{equation}
Consequently, given that either graph is RLCD, the other is RLCD exactly
when $N-2\nu-1=0$ in the field. The congruence alone does not establish
RLCD: for the five-cycle over $\mathbb F_2$ it holds, but adjacent
vertices have no common neighbour, so the corresponding entry of
$A(A+I)$ is one. Source Lemma~2 must be read as a compatibility statement
or corrected by adding the premise that one graph is RLCD.

For integral spectra, \Uref{ULCD.Rank} proves source Lemma~3: once $C$
is LCD, its dimension is the sum of the multiplicities of eigenvalues
nonzero modulo $p$. The proof uses the characteristic polynomial of the
actual integer matrix together with its LCD image/kernel decomposition;
it does not assume that reduction of a real eigenbasis is a basis over
$\mathbb F_p$.
The preliminary spectral observation on source page~5 also has its
precise elementary justification. For a symmetric integer matrix with
distinct integral eigenvalues $\lambda_i$, the integer polynomial
$\prod_i(X-\lambda_i)$ annihilates the matrix, since it does over
$\mathbb R$. If none reduce to zero, reduction gives invertibility.
If exactly one distinct $\lambda_i$ reduces to zero, the reduced
annihilator has only a simple factor $X$, so the generalized zero
eigenspace is the kernel. Its dimension is the multiplicity of zero in
the characteristic polynomial, namely the multiplicity of that
$\lambda_i$. With several eigenvalues reducing to zero, only the
general lower rank bound obtained from algebraic multiplicity is
automatic; the LCD premise supplies the exact formula used here.

Finally the decoding assertion quoted as source Result~1 is Massey's
nearest-neighbour observation. If $v=v_C+v_\perp$ is the orthogonal
splitting and $u\in C$ minimizes $\wt(v_\perp-u)$, then $v_C+u$ is a
nearest word in $C$, since translation by $v_C$ identifies the two
minimization problems. This is a statement about Hamming distance and
does not require an orthogonal sum formula for that distance. It gives
the source's decoder once the projections above are known.
